
\markright{সূচিপত্র}
\begin{center}
{\Huge\bfseries\color{chapterBlue}অধ্যায় ৫}\\[1.0mm]
{\LARGE\bfseries\color{darkTeal}প্রোগ্রামিং ভাষা}\\[0.9mm]
{\color{accentOrange}\rule{0.72\linewidth}{0.8pt}}
\end{center}
\vspace{0.7mm}
{\Large\bfseries\color{darkTeal}সূচিপত্র}\par
\vspace{0.55mm}

% দুই-কলামের এক-পৃষ্ঠার সূচিপত্র
{\footnotesize
\setlength{\tabcolsep}{2.6pt}%
\renewcommand{\arraystretch}{1.23}%
\begin{multicols}{2}
\setlength{\columnsep}{3.2mm}

\begin{tcolorbox}[
  enhanced,
  colback=white,
  colframe=lineGray,
  boxrule=0.45pt,
  arc=1mm,
  left=1.2mm,right=1.2mm,top=1mm,bottom=1mm
]
\arrayrulecolor{lineGray}
\begin{tabularx}{\linewidth}{@{}C{10.5mm}Y C{9mm}@{}}
\rowcolor{softBlue}\textbf{নং} & \textbf{বিষয়} & \textbf{পৃষ্ঠা}\\\hline
5.1 & প্রোগ্রামের ধারণা & \pageref{topic:1}\\\hline
5.2 & প্রোগ্রামিং ভাষার ধারণা & \pageref{topic:2}\\\hline
5.3 & মেশিন ভাষা (Machine Language) & \pageref{topic:3}\\\hline
5.4 & অ্যাসেম্বলি ভাষা (Assembly Language) & \pageref{topic:4}\\\hline
5.5 & মধ্যম স্তরের ভাষা (Middle Level Language) & \pageref{topic:5}\\\hline
5.6 & উচ্চস্তরের ভাষা (High Level Language) & \pageref{topic:6}\\\hline
5.7 & Java & \pageref{topic:7}\\\hline
5.8 & Oracle & \pageref{topic:8}\\\hline
5.9 & Python & \pageref{topic:9}\\\hline
5.10 & চতুর্থ প্রজন্মের ভাষা (4GL) & \pageref{topic:10}\\\hline
5.11 & অনুবাদক প্রোগ্রাম (Translator Program) & \pageref{topic:11}\\\hline
5.12 & কম্পাইলার (Compiler) & \pageref{topic:12}\\\hline
5.13 & অ্যাসেম্বলার (Assembler) & \pageref{topic:13}\\\hline
5.14 & ইন্টারপ্রেটার (Interpreter) & \pageref{topic:14}\\\hline
5.15 & প্রোগ্রামের সংগঠন (Program Organization) & \pageref{topic:15}\\\hline
5.16 & প্রোগ্রাম ডিজাইন মডেল (Program Design Model) & \pageref{topic:16}\\\hline
5.17 & মডুলার প্রোগ্রামিং (Modular Programming) & \pageref{topic:17}\\\hline
5.18 & প্রোগ্রাম ডিজাইনের কৌশল & \pageref{topic:18}\\\hline
5.19 & কম্পিউটেশনাল থিংকিং & \pageref{topic:19}\\\hline
5.20 & সমস্যাটি বিশ্লেষণ ও IPO মডেল & \pageref{topic:20}\\\hline
5.21 & প্রোগ্রাম ডেভেলপমেন্ট & \pageref{topic:21}\\\hline
5.22 & ড্রাই রান ও ট্রেস টেবিল & \pageref{topic:22}\\\hline
5.23 & অ্যালগরিদম (Algorithm) & \pageref{topic:23}\\\hline
5.24 & ফ্লোচার্ট (Flowchart) & \pageref{topic:24}\\\hline
5.25 & Python Programming Language-এর পরিচিতি & \pageref{topic:25}\\\hline
\end{tabularx}
\end{tcolorbox}
\columnbreak
\begin{tcolorbox}[
  enhanced,
  colback=white,
  colframe=lineGray,
  boxrule=0.45pt,
  arc=1mm,
  left=1.2mm,right=1.2mm,top=1mm,bottom=1mm
]
\arrayrulecolor{lineGray}
\begin{tabularx}{\linewidth}{@{}C{10.5mm}Y C{9mm}@{}}
\rowcolor{softBlue}\textbf{নং} & \textbf{বিষয়} & \textbf{পৃষ্ঠা}\\\hline
5.26 & Python Programming Language-এর বৈশিষ্ট্য & \pageref{topic:26}\\\hline
5.27 & Programming Environment & \pageref{topic:27}\\\hline
5.28 & Python Installation ও Setup & \pageref{topic:28}\\\hline
5.29 & Python-এ Code লেখার মাধ্যম & \pageref{topic:29}\\\hline
5.30 & Python Program-এর প্রাথমিক পরিচিতি & \pageref{topic:30}\\\hline
5.31 & Python Syntax-এর প্রাথমিক ধারণা & \pageref{topic:31}\\\hline
5.32 & Variable ও Data Type & \pageref{topic:32}\\\hline
5.33 & Python Variable & \pageref{topic:33}\\\hline
5.34 & Python Data Type & \pageref{topic:34}\\\hline
5.35 & Input ও Output & \pageref{topic:35}\\\hline
5.36 & Python Operator & \pageref{topic:36}\\\hline
5.37 & Conditional Logic Statement & \pageref{topic:37}\\\hline
5.38 & Loop — ধারণা ও প্রয়োজনীয়তা & \pageref{topic:38}\\\hline
5.39 & Python for Loop & \pageref{topic:39}\\\hline
5.40 & Python while Loop & \pageref{topic:40}\\\hline
5.41 & Loop Control Statement & \pageref{topic:41}\\\hline
5.42 & Loop Pattern & \pageref{topic:42}\\\hline
5.43 & Data Structure — তথ্যের সুশৃঙ্খল বিন্যাস & \pageref{topic:43}\\\hline
5.44 & Python-এ Function & \pageref{topic:44}\\\hline
5.45 & Function Contract, Scope, Local--Global ও File & \pageref{topic:45}\\\hline
5.46 & ব্যবহারিক Python Program অনুশীলন & \pageref{topic:46}\\\hline
5.47 & Exercise — Algorithm, Flowchart ও Python Program & \pageref{topic:47}\\\hline
5.48 & অধ্যায়ের সারসংক্ষেপ ও পরীক্ষাভিত্তিক প্রস্তুতি & \pageref{topic:48}\\\hline
5.49 & বহুনির্বাচনি প্রশ্ন (MCQ) & \pageref{topic:49}\\\hline
5.50 & সৃজনশীল প্রশ্ন (CQ) & \pageref{topic:50}\\\hline
-- & \textbf{সহায়ক গ্রন্থ ও তথ্যসূত্র} & \pageref{sources}\\\hline
\end{tabularx}
\end{tcolorbox}
\end{multicols}
}

% সূচিপত্র এক পৃষ্ঠায় শেষ; মূল অধ্যায় নতুন পৃষ্ঠা থেকে শুরু হবে।
\clearpage
\markright{অধ্যায় পরিচিতি}

\small

